% ═══════════════════════════════════════════════════════════════
% MT-08c: Minitest Preterito Indefinido (-ar Verben)
% Spanisch Klasse 7
% ═══════════════════════════════════════════════════════════════
% Dauer: 15 Minuten
% Punkte: 28 BE
% Inhalt: Signalwörter, Endungen, Konjugation -ar Verben
% ═══════════════════════════════════════════════════════════════

\documentclass[11pt, a4paper]{article}

\newif\ifswmodus
\swmodusfalse

\usepackage[utf8]{inputenc}
\usepackage[T1]{fontenc}
\usepackage[ngerman]{babel}
\usepackage{helvet}
\renewcommand{\familydefault}{\sfdefault}

\usepackage[margin=2cm]{geometry}
\usepackage{enumitem}
\usepackage{tabularx}
\usepackage{booktabs}

\usepackage{xcolor}
\usepackage{tcolorbox}
\tcbuselibrary{skins}

\usepackage{fancyhdr}
\usepackage{lastpage}
\usepackage{needspace}

% FARBEN
\definecolor{spanisch-color}{HTML}{c41e3a}
\definecolor{grau-color}{HTML}{888888}
\definecolor{hellgrau-color}{HTML}{f5f5f5}

\definecolor{sw-dark}{HTML}{000000}
\definecolor{sw-gray}{HTML}{666666}
\definecolor{sw-white}{HTML}{ffffff}

\ifswmodus
    \colorlet{spanisch}{sw-dark}
    \colorlet{grau}{sw-gray}
    \colorlet{hellgrau}{sw-white}
\else
    \colorlet{spanisch}{spanisch-color}
    \colorlet{grau}{grau-color}
    \colorlet{hellgrau}{hellgrau-color}
\fi

\newcommand{\fachfarbe}{spanisch}

\newcommand{\thema}{Minitest: Pretérito Indefinido (-ar)}
\newcommand{\fach}{Spanisch}
\newcommand{\klasse}{7}
\newcommand{\maxpunkte}{28}
\newcommand{\zeit}{15 Minuten}
\newcommand{\datum}{\today}

\pagestyle{fancy}
\fancyhf{}
\renewcommand{\headrulewidth}{0.4pt}
\renewcommand{\footrulewidth}{0.4pt}

\lhead{\fach{} -- Klasse \klasse}
\chead{\textbf{MT-08c}}
\rhead{\datum}

\lfoot{\zeit}
\cfoot{}
\rfoot{Seite \thepage\ von \pageref{LastPage}}

\newcommand{\punktebox}[1]{\hfill\fbox{\small \_/#1}}
\newcommand{\luecke}[1]{\rule{#1}{0.4pt}}

\newtcolorbox{infobox}{
    colback=hellgrau,
    colframe=\fachfarbe,
    boxrule=1pt,
    arc=0pt,
    left=8pt, right=8pt, top=8pt, bottom=8pt
}

\newenvironment{aufgabe}[1]{%
    \needspace{5\baselineskip}%
    \nopagebreak[4]%
    \vspace{10pt}
    \noindent\textbf{\textcolor{\fachfarbe}{Aufgabe #1}}
    \nopagebreak[4]%
    \vspace{5pt}
    \nopagebreak[4]%
    \begin{list}{}{\leftmargin=0pt}
    \item[]
}{%
    \end{list}
}

\newcommand{\es}[1]{\textcolor{\fachfarbe}{\textbf{#1}}}

\begin{document}

% ═══════════════════════════════════════════════════════════════
% KOPFBEREICH
% ═══════════════════════════════════════════════════════════════

\begin{center}
    {\Large\bfseries\textcolor{\fachfarbe}{\thema}}
\end{center}

\vspace{5pt}

\noindent
\begin{tabularx}{\textwidth}{|X|X|X|}
\hline
\textbf{Name:} \luecke{4cm} &
\textbf{Klasse:} \klasse &
\textbf{Datum:} \luecke{2cm} \\
\hline
\end{tabularx}

\vspace{10pt}

\begin{infobox}
\textbf{Zeit:} \zeit{} | \textbf{Punkte:} \maxpunkte{} BE | \textbf{Hilfsmittel:} keine \\
\textbf{Tipp:} Achte auf die Akzente: -é und -ó haben einen Akzent!
\end{infobox}

\vspace{15pt}

% ═══════════════════════════════════════════════════════════════
% AUFGABE 1: Signalwörter zuordnen (4 BE)
% ═══════════════════════════════════════════════════════════════

\begin{aufgabe}{1 -- Signalwörter zuordnen}
Verbinde das spanische Signalwort mit der deutschen Bedeutung. \punktebox{4}
\end{aufgabe}

\begin{center}
\begin{tabular}{rcl}
\es{ayer} & \luecke{2cm} & letztes Jahr \\[0.8em]
\es{el año pasado} & \luecke{2cm} & gestern \\[0.8em]
\es{la semana pasada} & \luecke{2cm} & letzten Sommer \\[0.8em]
\es{el verano pasado} & \luecke{2cm} & letzte Woche \\
\end{tabular}
\end{center}

\vspace{15pt}

% ═══════════════════════════════════════════════════════════════
% AUFGABE 2: Endungen zuordnen (6 BE)
% ═══════════════════════════════════════════════════════════════

\begin{aufgabe}{2 -- Endungen zuordnen}
Ergänze die richtige Endung für jede Person. \punktebox{6}
\end{aufgabe}

\noindent
\textbf{Endungen:} -é | -aste | -ó | -amos | -asteis | -aron

\vspace{0.5em}

\noindent
a) yo habl\luecke{2cm}

\vspace{0.8em}

\noindent
b) tú habl\luecke{2cm}

\vspace{0.8em}

\noindent
c) él/ella habl\luecke{2cm}

\vspace{0.8em}

\noindent
d) nosotros habl\luecke{2cm}

\vspace{0.8em}

\noindent
e) vosotros habl\luecke{2cm}

\vspace{0.8em}

\noindent
f) ellos habl\luecke{2cm}

\vspace{15pt}

% ═══════════════════════════════════════════════════════════════
% AUFGABE 3: Verbformen bilden (8 BE)
% ═══════════════════════════════════════════════════════════════

\begin{aufgabe}{3 -- Verbformen bilden}
Konjugiere die Verben im Pretérito Indefinido. \punktebox{8}
\end{aufgabe}

\noindent
a) yo / viajar: \luecke{4cm}

\vspace{0.8em}

\noindent
b) tú / nadar: \luecke{4cm}

\vspace{0.8em}

\noindent
c) él / visitar: \luecke{4cm}

\vspace{0.8em}

\noindent
d) nosotros / hablar: \luecke{4cm}

\vspace{15pt}

% ═══════════════════════════════════════════════════════════════
% AUFGABE 4: Lückentext (6 BE)
% ═══════════════════════════════════════════════════════════════

\begin{aufgabe}{4 -- Lückentext: Mis vacaciones}
Ergänze die Verben im Pretérito Indefinido. \punktebox{6}
\end{aufgabe}

\noindent
\textbf{El verano pasado} mi familia y yo \luecke{3cm} (viajar) a Mallorca.

\vspace{0.8em}

\noindent
Yo \luecke{3cm} (nadar) mucho en el mar.

\vspace{0.8em}

\noindent
Mi hermana \luecke{3cm} (tomar) el sol en la playa.

\vspace{15pt}

% ═══════════════════════════════════════════════════════════════
% AUFGABE 5: Übersetzen (4 BE)
% ═══════════════════════════════════════════════════════════════

\begin{aufgabe}{5 -- Übersetzen}
Übersetze die Sätze ins Spanische (Pretérito Indefinido). \punktebox{4}
\end{aufgabe}

\noindent
a) Ich reiste nach Spanien. $\rightarrow$ \luecke{6cm}

\vspace{0.8em}

\noindent
b) Wir besuchten Barcelona. $\rightarrow$ \luecke{6cm}

% ═══════════════════════════════════════════════════════════════
% PUNKTEÜBERSICHT
% ═══════════════════════════════════════════════════════════════

\vspace{25pt}
\hrule
\vspace{10pt}

\noindent
\begin{tabularx}{\textwidth}{|X|c|c|c|c|c||c|}
\hline
\textbf{Aufgabe} & 1 & 2 & 3 & 4 & 5 & \textbf{Gesamt} \\
\hline
\textbf{max. Punkte} & 4 & 6 & 8 & 6 & 4 & \textbf{28} \\
\hline
\textbf{erreicht} & & & & & & \\[0.8em]
\hline
\end{tabularx}

\vspace{10pt}

\begin{flushright}
\textbf{Note:} \luecke{2cm}
\end{flushright}

\end{document}
